% =============================================================================
%  1-Apertura.tex — Cuidados, alcance de la experiencia, objetivos y ruta
% =============================================================================


\section{Apertura: Cuidados, alcance de la experiencia, objetivos y ruta}
% ── Frame: Antes de empezar ──────────────────────────────────────────────────
\begin{frame}[shrink=6]{Antes de empezar}
  \setbeamerfont{block title}{size=\small}
  \noindent{\footnotesize Esta sesión incluye dinámicas con temas sensibles y no aptas para todo
  público; así como luces intensas, sonido fuerte y situaciones incómodas. Por lo
  que se recomienda leer los siguientes puntos antes de continuar:}

  \vspace{-0.05cm}
  \begin{columns}[T,totalwidth=\textwidth]
    \begin{column}{0.485\textwidth}
      \begin{block}{1. La participación es voluntaria}
        \footnotesize Puedes observar y no participar en las dinámicas si así lo deseas
        o bien, marcar tiempo fuera levantando la tarjeta roja proporcionada.
      \end{block}
      % \pause
      \begin{block}{3. Derecho de salida}
        \footnotesize Puedes salir del aula o retirarte de la actividad en cualquier
        momento; no necesitas dar explicaciones.
      \end{block}
    \end{column}
    \begin{column}{0.485\textwidth}
      \begin{block}{2. Sensibilidad a la luz y al sonido}
        \footnotesize Si tienes fotosensibilidad, epilepsia o alguna condición que pueda
        verse afectada por luces y sonidos fuertes, se recomienda no participar
        en las dinámicas.
      \end{block}
      % \pause
      \begin{block}{4. Aviso de ignorancia}
        \footnotesize En este grupo puede haber personas neurodivergentes, con o sin
        diagnóstico oficial. Para ellas, la sobrecarga no es una simulación; es
        una realidad.
      \end{block}
    \end{column}
  \end{columns}
\end{frame}

% ── Frame: Lo que no se puede hacer ──────────────────────────────────────────
\begin{frame}{Lo que no se puede hacer}
  {\color{ColorAcento}\rule{\textwidth}{0.8pt}}

  \vspace{0.15cm}
  {\normalsize\itshape\color{ColorPrincipal}
    «Simular discapacidades promueve el malestar y no mejora las actitudes hacia
    las personas con discapacidad»\par}
  \vspace{0.05cm}
  {\raggedleft\footnotesize\color{gray!70}
    --- \citet[p.~324]{nario2017}\par}

  \vspace{0.2cm}
  {\color{ColorAcento}\rule{0.6\textwidth}{0.5pt}}

  \vspace{0.15cm}
  {\normalsize\itshape\color{ColorPrincipal}
    «Nadie podrá llevar por encima de su corazón a nadie, ni hacerle mal en su
    persona, aunque piense y diga diferente.»\par}
  \vspace{0.05cm}
  {\raggedleft\footnotesize\color{gray!70}
    --- \citet{garzon1997discurso}\par}

  \vspace{0.15cm}
  {\color{ColorAcento}\rule{\textwidth}{0.8pt}}

  \vspace{0.15cm}
  {\small Por eso esta sesión no pretende que sepamos cómo es ser o padecer
  autismo. Busca visibilizar el papel del entorno y la hipersensibilidad: lo que
  el aula, la escuela, los compañeros y la sociedad hacen con quienes no encajan
  en la norma.\par}
\end{frame}

% ── Frame: Objetivos ─────────────────────────────────────────────────────────
\begin{frame}{Objetivos de la presentación}
  \begin{block}{Palabras clave}
    Vivir, comprender y empatizar.
  \end{block}
  \vspace{0.4cm}
  \begin{columns}[T]
    \begin{column}{0.32\textwidth}
      \centering
      {\color{ColorPrincipal}\Large\textbf{Vivir}}\\[0.1cm]
      {\small Aproximar la experiencia sensorial y cognitiva del autismo.}
    \end{column}
    \begin{column}{0.32\textwidth}
      \centering
      {\color{ColorPrincipal}\Large\textbf{Comprender}}\\[0.1cm]
      {\small Distinguir masking, hiperfoco, sobrecarga sensorial y sobrecarga cognitiva.}
    \end{column}
    \begin{column}{0.32\textwidth}
      \centering
      {\color{ColorPrincipal}\Large\textbf{Empatizar}}\\[0.1cm]
      {\small Conectar con la experiencia emocional y social del autismo.}
    \end{column}
  \end{columns}
\end{frame}

% ── Frame: Ruta de la sesión ─────────────────────────────────────────────────
\begin{frame}{Ruta de la sesión}
  \vfill
  \begin{enumerate}
    \setlength{\itemsep}{0.5cm}
    \item \textbf{Dinámica 1: El peso de la máscara}
    % \pause
    \item \textbf{Conceptos: Entendiendo la teoría a través del cansancio}
    % \pause
    \item \textbf{Dinámica 2-3-4: Un día normal en la escuela}
    % \pause
    \item \textbf{Tiempo para reflexión y preguntas}
  \end{enumerate}
  \vfill
\end{frame}
